%! Observational Studies — Unit 1, Lesson 1.5 — Homework (Answer Key)
\documentclass[10pt]{article}
\usepackage{saar-article}
\usepackage{saar-key}   % pulls in -boxes; do NOT also load -boxes
\newcommand{\TallMath}[1]{$\displaystyle #1\rule[-1.4em]{0pt}{3.2em}$}

\begin{document}

\pageheader{Unit 1, Lesson 1.5}{Homework --- Answer Key}
% NAMESTRIP: no \namedateperiod here — mirrors the blank. Teacher-only prose goes
% in the lesson plan as \begin{teachernote}[Homework], never in a key.

\vspace{0.15in}

\boxguard[24]
\begin{scenariobox}[Millbrook Public Library --- Does the Summer Reading Program Work?]{cerulean}
\small
Millbrook Public Library has $1{,}200$ card holders. Every summer it runs a
free Summer Reading Program that any card holder may sign up for. In September
the director draws a random sample of $240$ card holders from the full
membership list and reaches every one of them by phone.

\vspace{2pt}
Of the $240$, \textbf{$90$ had signed up for the program} and \textbf{$150$ had
not}. Of the $90$ who signed up, $72$ read at least $10$ books over the summer.
Of the $150$ who did not, $48$ did.
\end{scenariobox}

\vspace{0.08in}

\boxguard
\begin{notesbox}{Practice}
\small
\begin{enumerate}[leftmargin=1.6em, itemsep=9pt, topsep=3pt]

  \item The director did not put anybody into the program --- card holders
        signed themselves up.

        \textbf{(a)} Is this an observational study? \ans{yes}

        \textbf{(b)} Who decided which card holders were in each group?
        \ans{the card holders themselves}

  \item Compute the reading rate for each group.

        \textbf{(a)} What percent of the $90$ who signed up read at least $10$
        books?

        \begin{work}
          \dfrac{72}{90} &= 0.80 \\
                         &= 80\%
        \end{work}

        \textbf{(b)} What percent of the $150$ who did not sign up read at least
        $10$ books?

        \begin{work}
          \dfrac{48}{150} &= 0.32 \\
                          &= 32\%
        \end{work}

        \textbf{(c)} The program group is higher by \ans{$48$} percentage
        points.

  \item Now use the whole sample.

        \textbf{(a)} How many of the $240$ card holders read at least $10$
        books, and what percent is that?

        \begin{work}
          72 + 48 &= 120 \text{ card holders} \\
          \dfrac{120}{240} &= 0.50 = 50\%
        \end{work}

        \textbf{(b)} About how many of all $1{,}200$ card holders does that
        predict?

        \begin{work}
          0.50 \times 1200 &= 600 \text{ card holders}
        \end{work}

  \item Name the two variables, then write the finding the safe way.

        Explanatory: \ans{signing up for the program} \hfill Response: \ans{reading $10$ or more books}

        Write one sentence reporting the result as an \emph{association}, using
        both percents and no causal words.
        \ansline{Among Millbrook card holders, $80\%$ of those who signed up read at least $10$ books,}
        \ansline{against $32\%$ of those who did not --- signing up is associated with more reading.}

  \item The director's newsletter says: \emph{``The Summer Reading Program more
        than doubles how much our members read.''} Give \emph{two} explanations
        for the $48$-point gap other than the program itself.
        \ansline{(1) The arrow runs backwards: people who already read a lot are the ones who sign up.}
        \ansline{(2) A lurking variable --- families with free time both sign up and read more anyway.}

  \tcbbreak   % keep item 6's prompt with its table — mirrored in homework_key
  \item Each row below reports a real association somebody found by watching.
        Name one lurking variable that could explain it instead.

        \begin{center}
        \renewcommand{\arraystretch}{1.45}
        \begin{tabularx}{0.97\linewidth}{@{} X p{3.7cm} @{}}
        \toprule
        \textbf{What was observed} & \textbf{A lurking variable} \\
        \midrule
        On days when more ice cream is sold, more people get sunburned.
          & \ans{hot, sunny weather} \\
        Children with larger shoe sizes read at a higher level.
          & \ans{the child's age} \\
        Card holders who use the online catalog return books on time more often.
          & \ans{how organized the person is} \\
        Students who own a graphing calculator score higher on the state test.
          & \ans{the courses they are in} \\
        \bottomrule
        \end{tabularx}
        \end{center}

  \item The director checks one more record: how much each card holder read the
        \emph{previous} summer, before the program existed. Of the $90$ who
        signed up, $63$ had already read $10$ or more books. Of the $150$ who did
        not sign up, $30$ had.

        \textbf{(a)} Signed up: \ans{$70\%$} \qquad Did not sign up: \ans{$20\%$}

        \textbf{(b)} What does that say about who signs up, and does it make the
        program look better or worse than it is?
        \ansline{Heavy readers sign up: $70\%$ of them already read $10$ books, against $20\%$ of the}
        \ansline{others. The group started far ahead, so the program looks better than it really is.}

  \item The sample was drawn at random from the full card-holder list, every
        person answered, and every count is correct --- so there is no bias in
        this study.

        Explain why it \emph{still} cannot show that the program caused the
        extra reading.
        \ansline{A fair sample tells you the two groups really do differ; it does not tell you why. The}
        \ansline{card holders sorted themselves, so no random draw can rule out the other explanations.}

\end{enumerate}
\end{notesbox}

\vspace{0.08in}

\boxguard
\begin{extensionbox}
\small
Find a real headline that claims one thing causes another --- a news story, an
ad, a social media post. Copy it down exactly, then:

\begin{enumerate}[leftmargin=1.6em, itemsep=3pt, topsep=3pt]
  \item Name the explanatory variable and the response variable.
  \item Say whether the evidence behind it came from watching or from assigning.
  \item Name one lurking variable that could explain the association instead.
  \item Rewrite the headline so it claims only what the study can support.
\end{enumerate}
\end{extensionbox}

\vspace{0.08in}

\begin{remindbox}
\small
\textbf{Next class (Lesson 1.6):} an observational study cannot show causation
because the groups build themselves --- so what if we build them instead? Next
class is \textbf{experimental design}: assigning subjects to treatment and
control groups on purpose, and the four principles (comparison, randomization,
replication, and control) that make an experiment worth trusting.
\end{remindbox}

\end{document}
